\section{Data Construction and Coverage}\label{app:data}

\subsection{Crosswalk from O*NET--SOC to UK SOC 2020}

The exposure scores are indexed by O*NET--SOC occupation and the employment weights by UK SOC 2020, so the two are bridged by a published concordance that maps O*NET--SOC 2019 codes to UK SOC 2020 unit groups. I work at the three-digit minor-group level. This trades a degree of occupational resolution for cell sizes large enough to support reliable regional employment shares once survey waves are pooled, which is the binding constraint for Scotland, where finer cells thin out quickly. The coarser working level also makes the region-invariance assumption of \cref{as:reginv} more comfortable, since averaging over a minor group smooths the fine distinctions along which any residual regional difference in task content would appear.

\paragraph{Aggregation rule.} The concordance is many-to-many, so an aggregation rule is required. Each UK SOC 2020 unit group receives the mean of the scores of the O*NET occupations mapping into it, weighted by the concordance weights where the published file supplies them and unweighted otherwise. Unit groups then aggregate to the three-digit minor group by unweighted mean, since the public APS resolves employment only to the minor group and no domestic weight exists below it. A consequence of the rule is that the unweighted second step gives a thinly populated unit group the same voice as a dominant one. I therefore recompute the differential under two alternative rules, equal weights at both steps, and mapping-mass weights at the second step,
%, and a direct one-step weighted mean to the minor group, 
which \cref{tab:crosswalkrules} reports alongside the baseline. Ultimately, weighting the second step by UK unit-group employment would be the strongest approach, but it requires the secure-access APS and so is saved for future work granted such access.

\begin{table}[h]
\centering
\caption{Sensitivity of the differential to the crosswalk aggregation rule}
\label{tab:crosswalkrules}
\small
% Auto-generated by the analysis pipeline -- do not edit by hand.
% \input{} inside a table environment; requires \usepackage{booktabs}.
\begin{tabular}{l r r r}
\toprule
Aggregation rule & Scotland & rUK & Gap \\
\midrule
Baseline (weighted, then unweighted) & 0.6582 & 0.6669 & -0.008633 \\
Equal weights at both steps & 0.6582 & 0.6669 & -0.008633 \\
Mapping-mass weights at step 2 & 0.6595 & 0.6685 & -0.008988 \\
Direct one-step weighted mean & 0.6595 & 0.6685 & -0.008988 \\
\bottomrule
\end{tabular}

\begin{minipage}{0.92\textwidth}
\vspace{0.4em}\footnotesize
\textit{Notes:} Continuous composite index under the maximum operator, in index points. Step 1 maps O*NET--SOC occupations to UK SOC 2020 unit groups; step 2 maps unit groups to three-digit minor groups. The baseline row reproduces the headline gap of \cref{tab:headline} by construction.
\end{minipage}
\end{table}

\subsection{The differential at the four-digit level}

The three-digit working level compresses the compositional gap for the same reason coarse geographies compress spatial variation, as \cref{sec:differential} notes: any Scotland--rUK sorting across unit groups within a minor group is averaged away before the differential is formed. The public APS resolves occupation only to the minor group, so the comparison below requires the secure-access ASHE data for extension.


%\paragraph{Coverage.} \placeholder{share of UK employment mapping cleanly to a UK SOC 2020 minor group; number of unit groups requiring manual reconciliation, and the rule used. Printed by \texttt{02\_crosswalk.R}.}

\subsection{APS pooling and cell-count diagnostics}

I pool four APS waves, 2022--2025, applying the survey's person weights. Pooling is necessary because a single wave leaves too many Scottish minor-group cells sparsely populated for the employment shares to be precise, and the window is short enough that occupational structure is stable across it. The residual thinness is worth recording, because it is the source of the sampling error that dominates the reported interval. 23 of the 104 pooled Scottish occupation cells contain fewer than 100 respondents across the four waves, and the smallest contains 11. 
%The first-difference employment estimates weight by employment, which contains the damage but cannot repair it, and this is why I read the first-difference specification as uninformative rather than as evidence of no differential response.

%\subsection{The industry margin of the region-invariance assumption}
 
\paragraph{Two design features the bootstrap does not capture.} Consecutive APS annual datasets are not independent samples: the APS is assembled from overlapping Labour Force Survey waves, so a respondent can appear in adjacent annual datasets, and a person-level bootstrap of the pooled file treats those appearances as independent draws while also ignoring the survey's clustered design. Both simplifications understate the sampling variance of the shares, so the composed interval of \cref{tab:propagation} understates the sampling component and should not be read as an exact coverage statement. Separately, the pooled window coincides with the 2022--24 fall in LFS response rates and the associated reweighting, whose limitations for sub-national estimates the ONS itself has flagged; the bootstrap captures sampling variance but not differential non-response between Scotland and the rUK, which would bias the shares rather than widen their interval. The employer-reported ASHE dataset, available under a Data Access Agreement, is the natural external check on the pooled occupational shares, and I again leave that comparison to the secure-access data extension.

\Cref{as:reginv} fails on the industry margin exactly where the within-occupation industry mix differs across the border in a way correlated with exposure, and that much is checkable on the APS, which records industry alongside occupation. For each minor group $k$ I compute the total-variation distance between the Scottish and rUK industry distributions within the occupation, i.e.\ the share of within-occupation employment that would have to move between industries to equalise the two distributions,
\begin{equation*}
\mathrm{TV}_k = \tfrac{1}{2} \sum_i \bigl| \Pr(i \mid k, \Scot)
- \Pr(i \mid k, \rUK) \bigr|,
\end{equation*}
and summarise it in \cref{tab:industryinvariance}. ... ?
%Divergence is modest across most of the distribution, with a median TV of 0.14 and a 90th percentile of 0.26 across the 99 adequately populated minor groups. The Scotland-employment-weighted mean $B = \sum_k \sigma_{k,\Scot} \mathrm{TV}_k$ bounds the induced bias: if within-occupation exposure varies across industries by at most $\delta$ index points, the resulting error in $\dEbar$ cannot exceed $\delta B$. The bound is a sufficient condition and makes no claim about the true magnitude of $\delta$, which the design cannot observe. Read against the headline gap, the bound is informative only for modest $\delta$: with $B = 0.114$, a within-occupation industry spread of five index points permits an error of 0.57 points against a gap of $-0.86$. Two correlations speak to where the strain concentrates, and they pull in opposite directions. TV is mildly negatively correlated with the absolute gap contributions ($-0.19$), so the divergence does not concentrate in the cells that generate the differential. It is mildly \emph{positively} correlated with exposure itself ($+0.25$): occupations that are more exposed also tend to have somewhat more cross-border industry divergence. This does not on its own imply bias, since the bound depends on the product of TV with employment-share divergence rather than on exposure alone, and that product is what the negative correlation with gap contributions already speaks to; but it means the reassurance from the first correlation should not be read as unconditional, and I report both rather than the more flattering one alone. The adequate-cells threshold (thirty respondents per region) is looser than the roughly hundred-respondent standard used to flag thin cells elsewhere in this appendix, so the median, 90th-percentile, and most-divergent-cell figures partly reflect sampling noise in small industry-occupation cells and should be read as indicative of the shape of the distribution rather than as precise per-cell estimates.
 
%\begin{table}[h]
%\centering
%\caption{Within-occupation industry divergence and the implied bias bound}
%\label{tab:industryinvariance}
%\small
%% Auto-generated by the analysis pipeline -- do not edit by hand.
% \input{} inside a table environment; requires \usepackage{booktabs}.
\begin{tabular}{l r}
\toprule
Statistic & Value \\
\midrule
Minor groups with adequate cells & 99 \\
Share of Scottish employment covered & 0.9966 \\
Median TV & 0.1404 \\
90th percentile TV & 0.2606 \\
Bias bound multiplier B & 0.1143 \\
Implied bound on gap, delta = 5 index points & 0.5713 \\
Correlation of TV with exposure & 0.2474 \\
Correlation of TV with |gap contribution| & -0.1913 \\
\bottomrule
\end{tabular}

%\begin{minipage}{0.92\textwidth}
%\vspace{0.4em}\footnotesize
%\textit{Notes:} Computed on minor groups with at least thirty respondents in each region across the pooled 2022--2025 waves; coverage is reported as the share of Scottish employment those cells carry. TV is bounded in $[0,1]$, with zero indicating identical within-occupation industry composition on both sides of the border. $B$ sums $\sigma_{k,\Scot}\mathrm{TV}_k$ over the covered cells without renormalising, so it implicitly treats the uncovered 0.34\% of Scottish employment as contributing no divergence; the effect is negligible at this coverage level. The thirty-respondent floor is looser than the roughly hundred-respondent standard used elsewhere in this appendix, so the median and 90th-percentile rows should be read as indicative rather than precise. The final two rows test whether the strain on \cref{as:reginv} concentrates where the gap is generated.
%\end{minipage}
%\end{table}
 
\begin{table}[h]
\centering
\caption{Industry-mix and within-industry contributions by SIC section}
\label{tab:industrysections}
\small
% Auto-generated by the analysis pipeline -- do not edit by hand.
% \input{} inside a table environment; requires \usepackage{booktabs}.
\begin{tabular}{l l r r r r}
\toprule
Section & Industry & Share gap (Scot - rUK, pp) & Within-industry exposure gap & Industry-mix contribution & Within-industry contribution \\
\midrule
U & Extraterritorial bodies & -0.0004503 & -4.656 & -0.00002685 & -0.009024 \\
T & Household employers & -0.0504 & -2.141 & 0.008661 & -0.002386 \\
E & Water \textbackslash{}\& waste & 0.1663 & -1.939 & -0.002096 & -0.01772 \\
R & Arts \textbackslash{}\& recreation & 0.1401 & -1.372 & -0.001102 & -0.03904 \\
Q & Health \textbackslash{}\& social work & 1.364 & -1.315 & -0.02811 & -0.1998 \\
M & Professional \textbackslash{}\& technical & -1.706 & -1.281 & -0.1485 & -0.09176 \\
I & Accommodation \textbackslash{}\& food & 0.9293 & -1.16 & -0.1014 & -0.06816 \\
C & Manufacturing & -1.726 & -0.9774 & 0.002333 & -0.06307 \\
S & Other services & 0.08662 & -0.9615 & -0.005363 & -0.02747 \\
G & Wholesale \textbackslash{}\& retail & -0.2212 & -0.9404 & -0.00434 & -0.09922 \\
H & Transport \textbackslash{}\& storage & -0.6101 & -0.9088 & 0.01466 & -0.0384 \\
N & Administrative \textbackslash{}\& support & -0.35 & -0.8257 & 0.01769 & -0.03341 \\
O & Public admin \textbackslash{}\& defence & 1.532 & -0.5018 & 0.08306 & -0.04856 \\
F & Construction & -0.4983 & -0.3656 & 0.04711 & -0.02206 \\
J & Information \textbackslash{}\& communication & -1.414 & -0.3334 & -0.1446 & -0.01231 \\
L & Real estate & -0.1885 & -0.1449 & -0.01189 & -0.001545 \\
K & Finance \textbackslash{}\& insurance & 0.173 & -0.08175 & 0.02152 & -0.003637 \\
P & Education & -0.5587 & 0.7672 & 0.02436 & 0.07808 \\
A & Agriculture, forestry \textbackslash{}\& fishing & 0.5335 & 0.8904 & -0.06986 & 0.0122 \\
B & Mining \textbackslash{}\& quarrying & 1.732 & 0.9175 & 0.05923 & 0.01771 \\
D & Energy supply & 0.6671 & 1.145 & 0.02925 & 0.01348 \\
\bottomrule
\end{tabular}

\begin{minipage}{0.92\textwidth}
\vspace{0.4em}\footnotesize
\textit{Notes:} Per-section terms summing to the totals of \cref{tab:industry}, ordering A (rUK within-industry exposure, Scottish share differences); ordering B is materially similar section-by-section and its aggregate is reported in \cref{tab:industry}. The industry-mix term is demeaned, $\Delta\sigma_i (\Ebar_{i,\rUK} - \Ebar_{\rUK})$, so a section's entry reflects how unusual its Scotland--rUK share gap is rather than the level of $\Ebar_{i,\rUK}$; undemeaned terms and both orderings in full are in the online replication material. The industry names are abbreviated from the official SIC 2007 section titles.
\end{minipage}
\end{table}